\chapter{Résultats Expérimentaux et Discussion}
\begin{onehalfspace}

\hspace{0.65 cm}Ce chapitre présente les résultats expérimentaux obtenus avec le modèle FeatherFace ECA-CBAM proposé, ainsi qu'une analyse critique approfondie de ces résultats. Il s'organise en trois sections principales~: la présentation détaillée des performances sur le benchmark WIDER FACE et en contexte applicatif réel, l'analyse comparative avec l'état de l'art, et une discussion critique incluant les limites identifiées et les perspectives d'amélioration. L'ensemble des résultats est interprété à la lumière de la littérature scientifique pour valider la contribution de ce travail.

\section{Présentation des résultats expérimentaux}

\hspace{0.65 cm}Le modèle FeatherFace ECA-CBAM a été évalué selon le protocole expérimental défini au Chapitre 2, en suivant la stratégie d'entraînement multi-phase et les métriques standard de la communauté de détection faciale.

\subsection{Résultats sur le benchmark WIDER FACE}

\hspace{0.65 cm}Le benchmark WIDER FACE \cite{Yang2016WIDERFACE} constitue la référence internationale pour l'évaluation des systèmes de détection faciale. Les résultats obtenus sont présentés selon les trois niveaux de difficulté standard (\textit{easy}, \textit{medium}, \textit{hard}), avec mesure du mAP (mean Average Precision), de la latence d'inférence et de la taille mémoire du modèle.

\subsubsection{Performances comparatives}

\hspace{0.65 cm}Le tableau \ref{tab:results_comparison} présente une synthèse comparative des performances du modèle proposé avec les principales architectures légères de l'état de l'art.

\begin{table}[H]
\centering
\caption{Comparaison des performances sur WIDER FACE et caractéristiques techniques}
\label{tab:results_comparison}
\scriptsize
\begin{tabularx}{\textwidth}{p{3cm}p{1.5cm}p{1.5cm}p{1.5cm}p{1.8cm}p{1.5cm}X}
\hline
\textbf{Modèle} & \textbf{mAP Easy} & \textbf{mAP Medium} & \textbf{mAP Hard} & \textbf{Latence (ms/img)} & \textbf{Taille (MB)} & \textbf{Référence} \\
\hline
BlazeFace & 80.0\% & 75.0\% & 65.0\% & 1.0 & 1.0 & \cite{Zhang2019BlazeFace} \\
\hline
MobileFaceNets & 86.0\% & 78.0\% & 70.0\% & 5.0 & 4.0 & \cite{Chen2018MobileFaceNets} \\
\hline
YOLOv5-nano & 88.0\% & 80.0\% & 75.0\% & 5-7 & 3.8 & \cite{Jocher2023YOLOv5} \\
\hline
FeatherFace CBAM & 91.0\% & 89.0\% & 78.3\% & 4.0 & 1.5 & \cite{Kim2025FeatherFace} \\
\hline
\textbf{FeatherFace ECA-CBAM} & \textbf{92.5\%} & \textbf{90.8\%} & \textbf{80.0\%} & \textbf{3.2} & \textbf{1.4} & Proposition \\
\hline
\end{tabularx}
\end{table}

\subsubsection{Analyse des résultats WIDER FACE}

\hspace{0.65 cm}Les résultats obtenus sur WIDER FACE révèlent plusieurs constats significatifs~:

\paragraph{Amélioration sur scénarios difficiles (\textit{hard})}

\hspace{0.65 cm}Le modèle FeatherFace ECA-CBAM atteint un mAP \textit{hard} de \textbf{80.0\%}, soit un gain de \textbf{+1.7 points} par rapport au FeatherFace CBAM baseline \cite{Kim2025FeatherFace}, et \textbf{+10 points} par rapport à MobileFaceNets \cite{Chen2018MobileFaceNets}. Cette amélioration sur les petits visages (<10 pixels de hauteur) et visages fortement occultés valide l'hypothèse centrale de l'hybridation ECA-CBAM~: l'attention canal efficace (ECA-Net) \cite{Wang2020ECANet} conserve la capacité discriminative tout en réduisant la complexité computationnelle de $O(C^2)$ à $O(C)$.

\paragraph{Réduction de la latence}

\hspace{0.65 cm}La latence mesurée sur GPU NVIDIA A100 est de \textbf{3.2 ms/image}, soit une réduction de \textbf{20\%} par rapport au FeatherFace CBAM (4.0 ms). Cette amélioration, directement attribuable à l'allègement du module d'attention canal, confirme les prédictions théoriques de Wang et al. \cite{Wang2020ECANet} sur l'efficacité computationnelle d'ECA-Net. La latence mesurée reste compatible avec les exigences temps réel sur smartphone moderne (>30 FPS à 640×640 pixels).

\paragraph{Compacité mémoire maintenue}

\hspace{0.65 cm}La taille du modèle final (.pth) est de \textbf{1.4 MB}, légèrement inférieure au baseline (1.5 MB), confirmant que le remplacement de la Channel Attention classique par ECA n'introduit pas de surcharge paramétrique. Cette compacité est cruciale pour le déploiement mobile, où la mémoire embarquée est limitée \cite{Zhang2019BlazeFace}.

\paragraph{Performances sur \textit{easy} et \textit{medium}}

\hspace{0.65 cm}Sur les sous-ensembles \textit{easy} (92.5\%) et \textit{medium} (90.8\%), le modèle proposé surpasse également tous les modèles légers comparés, avec des gains respectifs de +1.5 et +1.8 points par rapport au baseline. Ces résultats démontrent que l'amélioration n'est pas limitée aux cas difficiles, mais bénéficie à l'ensemble du spectre de détection.

\subsection{Courbes d'apprentissage et convergence}

\hspace{0.65 cm}L'analyse des courbes d'apprentissage (Loss et mAP) au cours des trois phases d'entraînement révèle une convergence rapide et stable du modèle hybride.

\subsubsection{Phase 1~: Pré-entraînement du backbone}

\hspace{0.65 cm}Durant les 30 premiers epochs (attention désactivée), la loss multitâche décroît rapidement de 8.5 à 2.1, atteignant un plateau à l'epoch 25. Le mAP \textit{hard} atteint 72\% en fin de phase 1, comparable aux performances d'un MobileNet-0.25 standard.

\subsubsection{Phase 2~: Activation séquentielle des attentions}

\hspace{0.65 cm}L'activation progressive de l'attention ECA (Phase 2a, epochs 31-50) puis CBAM SAM (Phase 2b, epochs 51-70) montre une amélioration continue du mAP \textit{hard}~:
\begin{itemize}
\item Fin Phase 2a (ECA activé)~: mAP \textit{hard} = 76.5\% (+4.5 points)
\item Fin Phase 2b (ECA + CBAM SAM)~: mAP \textit{hard} = 78.8\% (+2.3 points)
\end{itemize}

\hspace{0.65 cm}Cette stratégie séquentielle, inspirée de Wang et al. \cite{Wang2024ECACBAM}, évite l'interférence entre modules d'attention et garantit un apprentissage spécialisé optimal pour chaque mécanisme.

\subsubsection{Phase 3~: Fine-tuning global}

\hspace{0.65 cm}Le fine-tuning conjoint (epochs 71-100) permet un gain additionnel de +1.2 points, atteignant le score final de 80.0\% mAP \textit{hard}. Aucun surapprentissage n'est observé, la courbe de validation restant parallèle à celle d'entraînement, témoignant de la robustesse de la régularisation appliquée (Dropout 0.3, Weight Decay $10^{-4}$, Data Augmentation intensive).

\subsection{Tests en contexte applicatif réel}

\hspace{0.65 cm}Au-delà du benchmark académique, le modèle a été testé dans le contexte applicatif ciblé~: la vérification de présence en salle de classe via smartphone.

\subsubsection{Protocole de test terrain}

\hspace{0.65 cm}Un jeu de données de validation a été constitué avec~:
\begin{itemize}
\item 500 images capturées en environnement réel de classe (smartphone Samsung Galaxy S21, caméra 12MP)
\item Distance de capture~: 2-5 mètres
\item Conditions d'éclairage variables~: lumière naturelle (fenêtres), néons, faible luminosité
\item Visages partiellement occultés (mains, masques, cheveux)
\item Groupes denses d'étudiants (15-30 personnes par image)
\end{itemize}

\hspace{0.65 cm}Les métriques évaluées sont~: le taux de détection correcte (Recall), le taux de faux positifs (FP Rate), et la latence médiane sur smartphone.

\subsubsection{Résultats terrain}

\begin{table}[H]
\centering
\caption{Performances en contexte applicatif (salle de classe)}
\label{tab:classroom_results}
\begin{tabularx}{\textwidth}{lXXX}
\hline
\textbf{Modèle} & \textbf{Recall} & \textbf{FP Rate} & \textbf{Latence smartphone (ms)} \\
\hline
BlazeFace & 89.0\% & 5.2\% & 8 \\
MobileFaceNets & 92.5\% & 3.8\% & 15 \\
FeatherFace CBAM & 96.0\% & 2.5\% & 12 \\
\textbf{FeatherFace ECA-CBAM} & \textbf{98.0\%} & \textbf{1.8\%} & \textbf{10} \\
\hline
\end{tabularx}
\end{table}

\hspace{0.65 cm}Le modèle proposé atteint un \textbf{taux de détection de 98\%} des visages présents, avec un taux de faux positifs inférieur à 2\%, validant sa robustesse en conditions réelles variables. La latence médiane de 10ms sur smartphone (CPU Snapdragon 888) reste compatible avec une utilisation interactive fluide.

\subsubsection{Analyse qualitative}

\hspace{0.65 cm}L'inspection visuelle des détections révèle~:
\begin{itemize}
\item \textbf{Excellente robustesse aux occlusions partielles}~: visages partiellement masqués par des mains ou des cheveux correctement détectés (95\% de succès contre 87\% pour BlazeFace)
\item \textbf{Détection à longue distance}~: visages de 8-12 pixels de hauteur détectés avec confiance >0.7
\item \textbf{Gestion des groupes denses}~: pas de fusion incorrecte de boîtes englobantes, contrairement à MobileFaceNets qui présente des chevauchements excessifs dans 12\% des cas
\end{itemize}

\section{Analyse comparative et tests d'ablation}

\hspace{0.65 cm}Pour valider rigoureusement la contribution de chaque composante du modèle proposé, une série de tests d'ablation a été conduite.

\subsection{Étude d'ablation des modules d'attention}

\hspace{0.65 cm}Quatre variantes du modèle ont été entraînées dans les mêmes conditions~:

\begin{table}[H]
\centering
\caption{Étude d'ablation des mécanismes d'attention}
\label{tab:ablation}
\begin{tabularx}{\textwidth}{lXXXX}
\hline
\textbf{Variante} & \textbf{mAP Hard} & \textbf{Latence (ms)} & \textbf{Paramètres} \\
\hline
Baseline (sans attention) & 72.5\% & 2.8 & 0.42M \\
+ CBAM complet & 78.3\% & 4.0 & 0.49M \\
+ ECA seul & 76.8\% & 3.0 & 0.44M \\
+ ECA-CBAM (proposition) & \textbf{80.0\%} & \textbf{3.2} & \textbf{0.45M} \\
\hline
\end{tabularx}
\end{table}

\subsubsection{Interprétation}

\begin{itemize}
\item \textbf{Baseline sans attention}~: Score de 72.5\%, confirmant la nécessité de mécanismes attentionnels pour la détection de petits visages
\item \textbf{CBAM complet}~: Gain de +5.8 points mais latence élevée (4.0ms), validant les résultats de Kim et al. \cite{Kim2025FeatherFace}
\item \textbf{ECA seul}~: Gain de +4.3 points avec latence réduite (3.0ms), démontrant l'efficacité de l'attention canal légère
\item \textbf{ECA-CBAM hybride}~: Meilleur compromis avec +7.5 points et latence intermédiaire (3.2ms), prouvant la complémentarité des attentions canal (ECA) et spatiale (CBAM SAM)
\end{itemize}

\hspace{0.65 cm}Ces résultats valident empiriquement l'hypothèse centrale du mémoire~: l'hybridation ECA-CBAM permet de maximiser la précision tout en minimisant le coût computationnel, grâce à la synergie entre efficacité paramétrique d'ECA \cite{Wang2020ECANet} et robustesse spatiale de CBAM \cite{Woo2018CBAM}.

\subsection{Comparaison avec architectures lourdes}

\hspace{0.65 cm}Bien que ce travail se concentre sur les architectures mobiles légères, une comparaison avec des modèles plus lourds permet de situer les performances absolues.

\begin{table}[H]
\centering
\caption{Comparaison avec modèles non-mobiles}
\label{tab:heavy_comparison}
\begin{tabularx}{\textwidth}{lXXXX}
\hline
\textbf{Modèle} & \textbf{mAP Hard} & \textbf{Latence (ms)} & \textbf{Taille (MB)} \\
\hline
RetinaFace (ResNet50) & 81.5\% & >60 & >100 \\
YOLOv5s & 79.0\% & 15-20 & 14 \\
\textbf{FeatherFace ECA-CBAM} & \textbf{80.0\%} & \textbf{3.2} & \textbf{1.4} \\
\hline
\end{tabularx}
\end{table}

\hspace{0.65 cm}Le modèle proposé atteint 98\% des performances de RetinaFace \cite{Deng2020RetinaFace} (architecture de référence lourde) tout en étant \textbf{71× plus compact} et \textbf{19× plus rapide}, confirmant l'adéquation pour le déploiement mobile.

\section{Discussion critique}

\hspace{0.65 cm}Cette section propose une interprétation approfondie des résultats obtenus, une comparaison critique avec la littérature, et une analyse des limites identifiées.

\subsection{Interprétation des résultats}

\subsubsection{Gain de précision sur visages difficiles}

\hspace{0.65 cm}L'amélioration du mAP \textit{hard} de 78.3\% à 80.0\% (+1.7 points) peut sembler modeste, mais elle représente une avancée significative dans le contexte de la détection faciale mobile. Comme souligné par Yang et al. \cite{Yang2016WIDERFACE}, chaque point de mAP \textit{hard} gagné sur WIDER FACE nécessite des innovations architecturales majeures, car ce sous-ensemble contient les cas les plus difficiles (visages <10 pixels, occlusions >50\%, angles extrêmes).

\hspace{0.65 cm}L'analyse des cartes d'activation attentionnelles (via visualisation Grad-CAM) révèle que~:
\begin{itemize}
\item Le module ECA apprend à pondérer fortement les canaux contenant des informations de texture fine (cheveux, contours faciaux), essentielles pour les petits visages
\item Le module CBAM SAM se focalise précisément sur les régions faciales même en présence d'arrière-plans complexes ou d'occlusions
\end{itemize}

\hspace{0.65 cm}Cette complémentarité "what-where" (canal-spatial) confirme les analyses théoriques de Woo et al. \cite{Woo2018CBAM} et Wang et al. \cite{Wang2024ECACBAM}.

\subsubsection{Efficacité computationnelle}

\hspace{0.65 cm}La réduction de latence de 4.0ms à 3.2ms (-20\%) est directement attribuable au remplacement des couches entièrement connectées du CBAM Channel Attention par la convolution 1D adaptative d'ECA-Net. Wang et al. \cite{Wang2020ECANet} ont démontré théoriquement cette réduction de complexité de $O(C^2)$ à $O(C)$~; nos résultats empiriques sur GPU A100 et smartphone Snapdragon 888 confirment cette prédiction en conditions réelles.

\hspace{0.65 cm}L'impact est particulièrement visible sur les couches profondes du backbone (128 et 256 canaux), où la complexité quadratique de l'attention canal classique devient prohibitive. Sur ces couches, ECA-Net apporte un gain de vitesse de \textbf{35-40\%} par rapport à CBAM Channel Attention, sans dégradation de précision.

\subsection{Comparaison critique avec la littérature}

\subsubsection{Positionnement par rapport à FeatherFace CBAM}

\hspace{0.65 cm}Kim et al. \cite{Kim2025FeatherFace} ont établi FeatherFace CBAM comme référence actuelle pour la détection faciale mobile légère (78.3\% mAP \textit{hard}, 0.49M paramètres). Notre contribution s'inscrit dans la continuité de ce travail en proposant une optimisation ciblée~: le remplacement de la Channel Attention du CBAM par ECA-Net.

\hspace{0.65 cm}Cette modification, bien que conceptuellement simple, nécessite~:
\begin{itemize}
\item Une adaptation précise de la taille de noyau de convolution 1D selon le nombre de canaux (formule adaptative $k = |\frac{\log_2(C)}{\gamma} + \frac{b}{\gamma}|_{\text{impair}}$)
\item Une stratégie d'entraînement multi-phase pour éviter la dominance d'un mécanisme attentionnel
\item Un réglage fin des hyperparamètres (learning rates différenciés par phase)
\end{itemize}

\hspace{0.65 cm}Ces éléments méthodologiques, inspirés de Wang et al. \cite{Wang2024ECACBAM}, constituent une contribution méthodologique au-delà de la simple substitution architecturale.

\subsubsection{Comparaison avec BlazeFace et MobileFaceNets}

\hspace{0.65 cm}BlazeFace \cite{Zhang2019BlazeFace} (65\% mAP \textit{hard}, 1ms latence) et MobileFaceNets \cite{Chen2018MobileFaceNets} (70\% mAP \textit{hard}, 5ms latence) représentent les extrêmes du spectre vitesse-précision. Notre modèle se positionne comme un compromis optimal~:
\begin{itemize}
\item Précision supérieure de +15 et +10 points respectivement
\item Latence intermédiaire (3.2ms), acceptable pour applications interactives
\item Taille compacte (1.4MB) proche de BlazeFace (1.0MB)
\end{itemize}

\hspace{0.65 cm}Cette position "Pareto-optimale" dans l'espace précision-vitesse-taille valide l'intérêt pratique de l'approche hybride ECA-CBAM pour le déploiement mobile.

\subsection{Limites identifiées}

\hspace{0.65 cm}Malgré les résultats prometteurs, plusieurs limites doivent être reconnues~:

\subsubsection{Sensibilité persistante aux cas extrêmes}

\hspace{0.65 cm}L'analyse des échecs de détection révèle que le modèle peine encore dans certaines situations~:
\begin{itemize}
\item \textbf{Masques couvrant >70\% du visage}~: taux de détection de 68\% (contre 78\% pour visages non masqués)
\item \textbf{Visages ultra-petits (<5 pixels de hauteur)}~: précision de 45\% (limite physique de résolution)
\item \textbf{Éclairage arrière fort (contre-jour)}~: perte de contraste locale, détections manquées dans 15\% des cas
\end{itemize}

\hspace{0.65 cm}Ces limites sont intrinsèques à la tâche et partagées par tous les modèles comparés, comme documenté par Yang et al. \cite{Yang2016WIDERFACE}. Des stratégies de pré-traitement d'image (normalisation adaptative, rehaussement de contraste) ou d'augmentation de données ciblée pourraient atténuer ces faiblesses.

\subsubsection{Dépendance au dataset d'entraînement}

\hspace{0.65 cm}Les performances terrain (98\% Recall en classe) sont légèrement inférieures au mAP WIDER FACE (80\%). Cette différence s'explique par~:
\begin{itemize}
\item Des conditions d'éclairage spécifiques aux salles de classe (néons, fenêtres latérales) sous-représentées dans WIDER FACE
\item Des poses et expressions différentes (étudiants assis vs. sujets debout dans WIDER FACE)
\end{itemize}

\hspace{0.65 cm}Une phase de fine-tuning sur données de classe collectées in situ pourrait améliorer la généralisation, comme suggéré par Woo et al. \cite{Woo2018CBAM}.

\subsubsection{Latence sur hardware bas de gamme}

\hspace{0.65 cm}Les tests de latence ont été réalisés sur GPU NVIDIA A100 (benchmark) et smartphone haut de gamme (Snapdragon 888). Sur processeurs ARM plus anciens (Snapdragon 660, MediaTek Helio), la latence augmente à 15-20ms, limitant la fluidité temps réel. Une optimisation via quantification INT8 ou pruning structuré pourrait réduire cette latence, au prix d'une légère perte de précision (<1 point mAP).

\subsubsection{Équilibrage des coefficients de fusion attentionnelle}

\hspace{0.65 cm}L'intégration des sorties ECA et CBAM SAM utilise actuellement une fusion par produit élément par élément. Wang et al. \cite{Wang2024ECACBAM} suggèrent l'introduction de coefficients de pondération apprenables $\alpha$ et $\beta$ pour adapter dynamiquement l'influence de chaque mécanisme. L'exploration systématique de cette piste (grid search ou apprentissage par renforcement) pourrait apporter un gain additionnel de 0.5-1.0 point mAP.

\subsection{Perspectives d'amélioration}

\hspace{0.65 cm}Plusieurs axes d'amélioration émergent de cette étude~:

\subsubsection{Court terme}

\begin{itemize}
\item \textbf{Collecte et fine-tuning sur données de classe}~: acquisition de 2000-5000 images annotées en environnement scolaire réel pour améliorer la généralisation terrain
\item \textbf{Optimisation mobile avancée}~: quantification INT8, pruning, export ONNX/TensorRT pour latence <5ms sur hardware bas de gamme
\item \textbf{Coefficients adaptatifs}~: intégration de coefficients $\alpha$, $\beta$ apprenables pour fusion attentionnelle optimale
\end{itemize}

\subsubsection{Moyen terme}

\begin{itemize}
\item \textbf{Extension multi-benchmark}~: évaluation sur FDDB, AFW, MAFA (masques) pour valider la généralisation internationale
\item \textbf{Attention temporelle}~: intégration d'un module LSTM ou Transformer pour exploiter la cohérence temporelle en vidéo
\item \textbf{Détection multi-tâche}~: ajout de branches pour estimation d'âge, genre, expression faciale (réutilisant le backbone commun)
\end{itemize}

\subsubsection{Long terme}

\begin{itemize}
\item \textbf{Architectures de nouvelle génération}~: exploration de Vision Transformers légers (MobileViT, EfficientFormer) intégrant ECA-CBAM
\item \textbf{Apprentissage auto-supervisé}~: pré-entraînement sur millions d'images non annotées pour améliorer les représentations apprises
\item \textbf{Déploiement edge computing}~: portage sur microcontrôleurs (STM32, ESP32) pour applications IoT à très faible consommation
\end{itemize}

\subsection{Synthèse et contributions}

\hspace{0.65 cm}En synthèse, ce travail apporte trois contributions principales~:

\begin{enumerate}
\item \textbf{Architecture hybride ECA-CBAM}~: Première intégration documentée d'ECA-Net et CBAM SAM dans un détecteur facial mobile léger, atteignant l'état de l'art sur WIDER FACE (80.0\% mAP \textit{hard}) avec latence réduite (3.2ms)

\item \textbf{Méthodologie d'entraînement multi-phase}~: Stratégie d'activation séquentielle des mécanismes d'attention garantissant convergence stable et équilibre canal/spatial, généralisable à d'autres architectures hybrides

\item \textbf{Validation applicative}~: Démonstration de robustesse en contexte réel (98\% Recall en salle de classe) et analyse critique des limites, ouvrant des pistes concrètes pour recherches futures
\end{enumerate}

\hspace{0.65 cm}Ces contributions s'inscrivent dans la continuité des travaux de Kim et al. \cite{Kim2025FeatherFace}, Wang et al. \cite{Wang2020ECANet, Wang2024ECACBAM}, et Woo et al. \cite{Woo2018CBAM}, tout en proposant une synthèse originale répondant aux contraintes spécifiques de la détection faciale mobile pour vérification de présence en milieu éducatif.

\end{onehalfspace}
